\chapter{Interoperabilidade entre marcas e modelos}
\label{cap:interoperabilidade}

\begin{objetivos}
Ao final deste capítulo você deve ser capaz de:
\begin{itemize}
  \item distinguir compatibilidade declarada, interoperabilidade e portabilidade;
  \item aplicar as estratégias que realmente reduzem dependência de fabricante;
  \item localizar, no projeto, os pontos em que a migração quebra;
  \item classificar o risco de uma migração antes de orçá-la;
  \item escrever exigências contratuais que preservem a liberdade futura da planta.
\end{itemize}
\end{objetivos}

\section{Três palavras que não são sinônimos}

No material de fornecedor, as três aparecem trocadas:

\begin{itemize}
  \item \textbf{compatibilidade declarada}: o fabricante afirma que o produto atende a uma norma. É
        uma declaração, não uma verificação independente --- e cabe perguntar a qual edição da norma
        ela se refere (Capítulo~\ref{cap:familia-iec});
  \item \textbf{interoperabilidade}: equipamentos diferentes conseguem trocar dados em operação.
        Depende de protocolo e de perfil de informação, não de linguagem de programação;
  \item \textbf{portabilidade}: o mesmo \emph{programa} roda em outro ambiente com esforço
        previsível. É a propriedade mais difícil das três, porque depende de tudo o que a norma
        \emph{não} padroniza.
\end{itemize}

Interoperabilidade resolve o problema de \emph{conversar}. Portabilidade resolve o problema de
\emph{migrar}. São problemas distintos e frequentemente confundidos em conversa de especificação.

\index{interoperabilidade}\index{portabilidade}

\section{A escada da portabilidade}

\begin{figure}[H]
  \centering
  \includegraphics[width=0.50\linewidth]{20-portabilidade.pdf}
  \caption{O que é portável e o que cada ambiente acrescenta por conta própria.}
  \label{fig:portabilidade}
\end{figure}

A Figura~\ref{fig:portabilidade} mostra que a portabilidade não é uma propriedade binária: ela
diminui à medida que se desce na pilha. No topo, linguagens, POUs e tipos de dados são padronizados e
viajam bem. Nos degraus seguintes, cada ambiente acrescenta o seu --- e é aí que o custo aparece.

O Capítulo~\ref{cap:iec-versus-legado} já separou os problemas de linguagem dos problemas de
estrutura e mostrou que os primeiros são mecânicos e os segundos não. Este capítulo trata os
segundos.

\section{Estratégias que funcionam}

\begin{enumerate}
  \item \textbf{escrever conforme a norma, sempre que houver equivalente padronizado}. É a decisão de
        maior efeito e de menor custo, tomada no dia da escrita --- não no dia da migração
        (Tabela~\ref{tab:recomendacoes-portabilidade});
  \item \textbf{usar biblioteca padronizada}, e não blocos proprietários, para as funções que existem
        em versão normatizada. Bibliotecas setoriais de blocos de função existem justamente para
        isso (Capítulo~\ref{cap:familia-iec});
  \item \textbf{isolar o que é específico}: endereçamento, diagnóstico de rack, comunicação e
        qualquer função de sistema devem ficar concentrados em poucos blocos, com nome e comentário
        explícitos;
  \item \textbf{preferir texto estruturado} para lógica complexa. Ladder é excelente para lógica
        combinacional e para quem lê em campo, mas algoritmos longos ficam mais legíveis e mais
        facilmente transportáveis em ST;
  \item \textbf{usar protocolo aberto como ponte} quando o código não é portável. Nem sempre vale
        reescrever o programa: muitas vezes a integração se resolve mantendo os dois controladores e
        ligando-os por OPC UA
        (Capítulo~\ref{cap:iiot}) ou Modbus TCP (Capítulo~\ref{cap:redes-scada});
  \item \textbf{usar o formato de intercâmbio} previsto na família da norma para exportar e importar
        projetos entre ambientes, quando ambos o implementarem. Ele transporta a estrutura, não as
        bibliotecas proprietárias.
\end{enumerate}

\begin{atencao}
As estratégias 1 a 4 custam quase nada se aplicadas desde o início e custam muito quando aplicadas
depois. Reescrita de programa em operação consome parada de produção, teste e revalidação,
comissionamento e risco de introduzir defeito em código que já funcionava. Quem decide sobre
arquitetura de software conversa com quem decide sobre orçamento de manutenção --- e raramente são a
mesma pessoa.
\end{atencao}

\section{Onde a migração realmente quebra}

\begin{table}[H]
  \centering
  \caption{Pontos de ruptura em uma migração de fabricante e o que fazer com cada um.}
  \label{tab:rupturas-migracao}
  \begin{tabular}{p{3.4cm} p{5.2cm} p{5.7cm}}
    \toprule
    Ponto de ruptura & Por que quebra & O que reduz o impacto \\
    \midrule
    Endereçamento físico de E/S
      & Cada família organiza e nomeia os pontos de um jeito. A norma não padroniza endereçamento
      & Nomear variáveis pelo processo, nunca pelo endereço; manter tabela de alocação
        (Capítulo~\ref{cap:documentacao}) \\
    \addlinespace
    Blocos de sistema
      & Leitura de estado do módulo, diagnóstico de rack e tratamento de falha são próprios de cada
        ambiente
      & Concentrar em um bloco de adaptação, documentado; testar esse bloco primeiro na migração \\
    \addlinespace
    Temporização e contagem
      & Parâmetros, base de tempo e comportamento de acumulação variam entre ambientes
      & Declarar explicitamente base de tempo e comportamento esperado; testar com cronômetro, não
        por leitura de código \\
    \addlinespace
    Comunicação entre controladores
      & Blocos de troca de dados proprietários não têm equivalente direto
      & Padronizar em protocolo aberto desde o projeto, ou em camada de gateway \\
    \addlinespace
    HMI e SCADA
      & Telas, animações e \emph{scripts} ficam presos ao produto e à versão
      & Separar o endereçamento da tela do endereço físico, com camada de tags intermediária \\
    \addlinespace
    Formato de projeto e bibliotecas
      & Arquivo, versionamento e bibliotecas são do fabricante
      & Exigir entrega do projeto-fonte e exportação em formato aberto; manter o fonte em versionamento
        próprio \\
    \addlinespace
    Peças e ciclo de vida
      & Migração forçada por falta de reposição é a mais cara de todas
      & Acompanhar avisos de descontinuação; planejar substituição antes de a parada ser obrigatória \\
    \bottomrule
  \end{tabular}
\end{table}

\section{Classificando o risco da migração}

Nem toda migração é igual. A Tabela~\ref{tab:risco-migracao} dá uma escala de trabalho que ajuda a
dimensionar esforço antes de orçar.

\begin{table}[H]
  \centering
  \caption{Cenários de migração e esforço típico.}
  \label{tab:risco-migracao}
  \begin{tabular}{p{5.0cm} p{2.4cm} p{6.8cm}}
    \toprule
    Cenário & Risco & Por quê \\
    \midrule
    Troca de modelo dentro da mesma família, com o mesmo ambiente de programação
      & Baixo
      & Endereçamento, blocos e formato de projeto permanecem; o esforço é de conferência e teste
        de compatibilidade de \emph{firmware} \\
    \addlinespace
    Troca de família com o mesmo fabricante e o mesmo ambiente
      & Médio
      & A estrutura de projeto e a biblioteca mudam; a linguagem, em geral, não \\
    \addlinespace
    Troca de fabricante, mantendo o projeto funcional e a arquitetura de rede
      & Alto
      & Reescrever aplicativo, refazer endereçamento, revalidar comunicação e reescrever a interface
        com a supervisão \\
    \addlinespace
    Troca de fabricante com substituição de HMI, supervisão e rede de campo
      & Muito alto
      & É um projeto novo, com comissionamento novo; deve ser orçado como tal, e não como
        ``atualização de controlador'' \\
    \bottomrule
  \end{tabular}
\end{table}

Não existe, nesta tabela, promessa de compatibilidade entre marcas. Existe uma escala de esforço
verificável no campo: quanto mais camadas mudam ao mesmo tempo, mais caro fica --- e é justamente por
isso que quem migra ganha ao migrar por partes.

\index{migração de controlador}

\section{O que exigir na especificação e no contrato}

Toda liberdade futura da planta é decidida no momento da compra. Do que foi discutido até aqui,
quatro exigências resolvem a maior parte do problema (ver também o Capítulo~\ref{cap:porte}):

\begin{enumerate}
  \item \textbf{entrega do projeto-fonte}, não apenas do arquivo executável carregado no controlador,
        com a versão do ambiente de programação identificada;
  \item \textbf{declaração de conformidade com edição e ano} da norma, e a informação de quais
        linguagens e bibliotecas foram efetivamente usadas;
  \item \textbf{lista de blocos proprietários} utilizados no aplicativo, com o que cada um faz ---
        essa lista é o orçamento da migração futura;
  \item \textbf{exportação em formato de intercâmbio} sempre que o ambiente oferecer, e a tabela de
        tags em formato aberto (planilha ou texto), nunca apenas dentro do projeto.
\end{enumerate}

\begin{pratica}
Escolha um sistema real com pelo menos um controlador, uma IHM e uma rede de campo. Produza um
relatório de portabilidade com:
(a) a lista de blocos proprietários usados no programa, com a função de cada um;
(b) a classificação do cenário de migração na Tabela~\ref{tab:risco-migracao};
(c) os três pontos de ruptura da Tabela~\ref{tab:rupturas-migracao} que exigiriam mais trabalho no
seu caso, justificando;
(d) as exigências contratuais que faltaram na compra original.
\end{pratica}

\begin{exercicios}
\begin{enumerate}[label=\alph*)]
  \item Explique a diferença entre interoperabilidade e portabilidade com um exemplo de cada.
  \item Por que a portabilidade diminui à medida que se desce na Figura~\ref{fig:portabilidade}?
  \item Um projeto usa apenas LD e blocos \texttt{TON} padronizados, mas lê o estado dos módulos por
        bloco proprietário. Classifique o risco de migração e justifique.
  \item Cite dois casos em que manter dois controladores ligados por protocolo aberto é mais
        vantajoso que reescrever o programa.
  \item Por que nomear variáveis pelo processo, e não pelo endereço, sobrevive à migração?
  \item Explique o risco de temporização e diga por que a conferência precisa ser feita com
        cronômetro.
  \item Classifique na Tabela~\ref{tab:risco-migracao}: (i) troca de um controlador por outro do mesmo
        modelo, por fim de vida; (ii) substituição completa de um painel, com controlador, IHM e rede
        nova; (iii) troca de fabricante mantendo a IHM existente.
  \item Justifique, com base na Tabela~\ref{tab:rupturas-migracao}, por que a migração forçada por
        falta de peça é a mais cara.
  \item Escreva três cláusulas de especificação para a compra de um controlador, cada uma ligada a um
        ponto de ruptura.
  \item Um fornecedor afirma: ``nosso controlador é conforme IEC 61131-3, então seu programa roda em
        qualquer marca''. Escreva a resposta técnica em cinco linhas.
\end{enumerate}
\end{exercicios}
